\section{Introduction}

Let \(\rho\) be a non-Gaussian probability law and
\(\mu=\rho^{\otimes\N}\).  We ask which bounded invertible operators on
\(\ell^2\) preserve the measure class of \(\mu\) when their rows act on the
coordinate sequence.  Kakutani's theorem answers the corresponding question
for coordinatewise transformations, but general linear mixing destroys the
independence on which its product criterion relies.  Fixed finite output
marginals do not resolve the problem: they may all have positive densities
even when the full image law is singular.

For finite-variance coordinate laws satisfying the regularity assumptions
below, the answer has the same operator-ideal form as the
Feldman--H\'ajek criterion, but with a discrete non-Gaussian stabilizer.  The
image law is equivalent to \(\mu\) precisely when the mixing operator is a
Hilbert--Schmidt perturbation of an allowed signed coordinate permutation;
equivalently, the Hellinger energies of its finite output marginals are
uniformly bounded in the number of output coordinates.

Stable laws and coordinate densities with zeros show how tail and positivity-set
geometry alter this criterion in different ways.

\subsection{Main theorem}

Let \(X\sim\rho\).  Suppose \(\rho(dx)=f(x)\,dx\), put \(g=\sqrt f\), and assume
\begin{equation}
 \begin{gathered}
 f>0\quad\text{Lebesgue-almost everywhere},\qquad
 g\in W^{1,2}(\R),\qquad xg'\in L^2(\R),\\
 \E X=0,\qquad \E X^2=1,\qquad
 f\ne(2\pi)^{-1/2}e^{-x^2/2}.
 \end{gathered}
 \label{eq:main-assumptions}
\end{equation}
Derivatives are weak; these hypotheses give finite location and scale Fisher
information without requiring moments above order two.

Put \(H=\ell^2(\N;\R)\), let \((e_i)_{i\ge1}\) be its coordinate basis, and
write \(\mathcal S_2(H)\) for the Hilbert--Schmidt ideal~\cite{simon2005trace}.
Define
\begin{equation}
 \begin{aligned}
 \mathcal P=\{Q:{}&(Qx)_i=\varepsilon_i x_{\pi(i)},\quad
       \pi\text{ a permutation of }\N,\\
       &\varepsilon_i\in\{-1,1\},\quad
       \Law(\varepsilon_iX)=\rho\text{ for every }i\}.
 \end{aligned}
\end{equation}
Thus \(\mathcal P\) is the full signed-permutation group when \(\rho\) is
symmetric and the ordinary permutation group otherwise.  For
\(T\in\mathcal B(H)\), let \(\Phi_T\) be the measurable row-series
realization constructed in
\cref{lem:measurable-row-series-realization}, and set
\(\nu_T=(\Phi_T)_\#\mu\).  A sample from \(\mu\) is not assumed to lie in
\(H\).

For probability measures \(\lambda,\nu\) and any common dominating measure
\(m\), write
\begin{equation}
 \Aff(\lambda,\nu)
 =\int\sqrt{\frac{d\lambda}{dm}\frac{d\nu}{dm}}\,dm,
 \qquad
 h(\lambda,\nu)^2=2-2\Aff(\lambda,\nu).
\end{equation}
Let \(\mu_n=\rho^{\otimes n}\), let \(\pi_n\) extract the first \(n\)
coordinates, and define
\begin{equation}
 E_n(T)=-2\log\Aff\bigl(\mu_n,(\pi_n)_\#\nu_T\bigr).
 \label{eq:projected-hellinger-energy}
\end{equation}
These energies use the first \(n\) output rows of \(T\), not its principal
\(n\times n\) corner.

\begin{theorem}[Bounded invertible linear mixing]
\label{thm:measure-equivalence}
Under \eqref{eq:main-assumptions}, every \(T\in\GL(H)\) satisfies
\begin{equation}
 \boxed{\begin{aligned}
 \nu_T\sim\mu
 &\quad\Longleftrightarrow\quad
 \exists Q\in\mathcal P:\ TQ^{-1}-I\in\mathcal S_2(H)\\
 &\quad\Longleftrightarrow\quad \sup_nE_n(T)<\infty.
 \end{aligned}}
\end{equation}
If these conditions fail, then \(\nu_T\perp\mu\) and
\(E_n(T)\uparrow\infty\).  Moreover, \(\nu_T=\mu\) if and only if
\(T\in\mathcal P\).
\end{theorem}

Bounded projected energy first detects a nonzero absolutely continuous
component; the equivalence--singularity dichotomy upgrades that overlap to equivalence.

\begin{corollary}[Moment-only stabilizer]
\label{cor:moment-only-stabilizer}
If \(\rho\) is non-Gaussian, centered, and has variance one, then the exact
linear stabilizer of \(\rho^{\otimes\N}\) is \(\mathcal P\).  No density or
Fisher information assumption is needed.
\end{corollary}

\begin{corollary}[Orthogonal transformations]
\label{cor:orthogonal-equivalence}
If \(f>0\) almost everywhere, \(\sqrt f\in W^{1,2}(\R)\), and \(f\) is
non-Gaussian, centered, and has variance one, then every orthogonal \(U\)
satisfies
\[
 \nu_U\sim\mu
 \quad\Longleftrightarrow\quad
 U-Q\in\mathcal S_2(H)\quad\text{for some }Q\in\mathcal P.
\]
Thus scale Fisher information is not needed for orthogonal mixing.
\end{corollary}

\subsection{Relation to earlier criteria}

Kakutani characterizes equivalence of products with equivalent factors by
summability of one-coordinate Hellinger losses~\cite{kakutani1948equivalence}; for translations, Shepp
relates this sum to square-root-density regularity
\cite{shepp1965translate}.  Euclidean and Lie-group perturbations identify
Fisher information as the local metric while retaining product structure
\cite{steele1986fisher,marques1992quasi}.  General output rows instead require
estimates sensitive to the dependence created by mixing.

For Gaussian measures, Cameron--Martin describes admissible translations
\cite{cameron1944transformations}, while Feldman--H\'ajek theory and likelihood
formulas describe equivalence under linear transformations
\cite{feldman1958equivalence,hajek1958property,shepp1966radon,bogachev1998gaussian}.
Ramer gives sufficient quasi-invariance and change-of-variables conditions for
nonlinear Gaussian transformations
\cite{ramer1974nonlinear}.  Every orthogonal operator is an exact symmetry of
the standard Gaussian product; in the non-Gaussian case the stabilizer is
discrete, and equivalence is determined by Hilbert--Schmidt proximity to that
stabilizer rather than by the covariance defect alone.

Shimomura's Theorem~3.4 shows that \(I-S^2\) is compact for a self-adjoint
admissible operator \(S\), under \(\ell^2\)-continuity,
\(\ell^2\)-translation quasi-invariance, and \(\ell^2\)-ergodicity; his
examples preclude a universal Hilbert--Schmidt criterion under those
assumptions alone~\cite{shimomura1977linear}.  Under
\eqref{eq:main-assumptions}, \cref{thm:measure-equivalence} gives the exact
Hilbert--Schmidt condition modulo \(\mathcal P\) for every bounded invertible
operator.  Other works construct non-Gaussian quasi-invariant measures on
linear spaces~\cite{feldman1961examples,takahashi1975quasi}; here we fix one
i.i.d. product and classify its admissible linear maps.

Classical independence-of-linear-forms theorems
\cite{kac1939characterization,darmois1953analyse,skitovich1954linear,kagan1973characterization}
characterize Gaussian laws through independence of suitable linear forms,
while independent component analysis identifies non-Gaussian mixing matrices
up to scaling and permutation~\cite{comon1994independent}.  Here the output
coordinates may be dependent; their finite marginals instead detect
square-summable deviation from an allowed signed coordinate permutation.

\subsection{Two boundary regimes}

For the standard symmetric \(\alpha\)-stable product, \(0<\alpha<2\),
\cref{thm:complete-criterion-symmetric-stable-coordinates} shows, within the
compatible stable row-series class, that equivalence holds precisely when an
allowed signed coordinate permutation reduces the operator to a matrix \(M=(m_{ij})\)
satisfying
\[
 \sum_i|m_{ii}-1|^2
 +\sum_j\left(\sum_{i\ne j}|m_{ij}|^2\right)^{\alpha/2}<\infty.
\]
The off-diagonal term is an \(\ell^\alpha(\ell^2)\) column cost, so positive
density and finite Fisher information do not imply Hilbert--Schmidt
sufficiency.

When a finite-variance density may vanish, Hilbert--Schmidt necessity persists,
but \cref{thm:linear-equivalence-group} additionally requires every row sum
of both \(T\) and \(T^{-1}\) to lie almost surely in \(\{f>0\}\).  For bounded
or one-sided coordinate laws,
\cref{thm:bounded-and-half-line-classification} reduces the equivalence group
to explicit signed-permutation or permutation classes without density or
Fisher assumptions.

\subsection{Proof strategy}

The finite-variance proof has three main steps.  First, non-Gaussianity is
encoded by the rotational Fourier defect
\[
 R(\xi,\eta)=\eta\phi'(\xi)\phi(\eta)
              -\xi\phi(\xi)\phi'(\eta),
 \qquad \phi(u)=\E e^{iuX}.
\]
The determinant of the two weighted Fourier jets
\((\xi\phi(\xi),\phi'(\xi))\) and
\((\eta\phi(\eta),\phi'(\eta))\) vanishes at every pair only for a Gaussian
law.  A finite set of frequencies then yields bounded tests for all matrix
directions.  Combined with the matrix Fisher operator, these tests give local
Hellinger coercivity in Frobenius norm with constants independent of the
dimension.

Second, the operator action must be realized on \(\R^{\N}\), since a product
sample need not belong to \(\ell^2\).  A common measurable carrier makes row
series composable and allows inverse maps and likelihood cocycles to be used
without incompatible null sets.  Translation ergodicity gives the
equivalence--singularity dichotomy and identifies bounded projected energy with
equivalence.

Third, absolute continuity initially gives only asymptotic matching of remote
rows with allowed signed coordinates.  If the squared matching errors were
not summable, tail-coordinate swaps would package a fixed amount of error into
a small Hilbert--Schmidt perturbation.  Local coercivity detects this
perturbation, while absolute continuity makes the same remote swap
asymptotically invisible.  The contradiction proves Hilbert--Schmidt
necessity.

We use \(A^*\) for the adjoint, \(P_n\) for coordinate projection on \(H\),
and \(\|\cdot\|_F\) for the Frobenius norm.
